\section{Fourier Series}\label{sec.fourier_series}

\recall the general solution to the heat equation \eqref{eqn.heat_eqn_sol}:
\[
u(x,t) = \sum_{n=1}^\infty b_n \sin\left(\frac{n \pi x}{L}\right) e^{-a^2 \left(\frac{n \pi}{L}\right)^2 t}
\]

The goal of this section is to build a method for finding the coefficients $b_n$ in the above solution, given an initial condition. To do this, we will need to introduce the concept of Fourier series, which allows us to express a function as an infinite sum of sines and cosines.

\subsection{Fourier Series}\label{subsec.fourier_series}
A Fourier series is a way to represent a periodic function as an infinite sum of sines and cosines. For a function $f(x)$ defined on the interval $[0, L]$, the Fourier series is given by:
\begin{equation*}\label{eqn.fourier_series}
f(x) = \frac{a_0}{2} + \sum_{n=1}^\infty \left( a_n \cos\left(\frac{n \pi x}{L}\right) + b_n \sin\left(\frac{n \pi x}{L}\right) \right)
\end{equation*}
where the coefficients $a_n$ and $b_n$ are given by:
\begin{equation*}\label{eqn.fourier_a_coeffs}
a_n = \frac{2}{L} \int_0^L f(x) \cos\left(\frac{n \pi x}{L}\right) dx, \quad n \geq 0
\end{equation*}
\begin{equation*}\label{eqn.fourier_b_coeffs}
b_n = \frac{2}{L} \int_0^L f(x) \sin\left(\frac{n \pi x}{L}\right) dx, \quad n \geq 1
\end{equation*}

So we want to find the coefficients, $b_n$, of the Fourier basis functions, $\sin\left(\frac{n \pi x}{L}\right)$, to represent $f$ as a sum of sines.
\subsubsection{Why are they a basis?}\label{subsubsec.fourier_basis}
To show this properly, we need to review some topics from linear algebra through the following definitions.

\begin{definition}\label{def.lin_indep}
    A set of functions $\{f_n\}$ is said to be \textbf{linearly independent} if the only way to write the zero function as a linear combination of the $f_n$ is to have all coefficients be zero. In other words, if $\sum_{n} c_n f_n(x) = 0$ for all $x$, then $c_n = 0$ for all $n$.
\end{definition}

\begin{definition}\label{def.basis}
    A set of functions $\{f_n\}$ is a \textbf{basis} for a function space if every function in the space can be written as a linear combination of the $f_n$. In other words, for any function $g$ in the space, there exist coefficients $c_n$ such that $g(x) = \sum_{n} c_n f_n(x)$ for all $x$.
\end{definition}

\begin{definition}\label{def.inner_product}
    An \textbf{inner product} on a function space is the same as an inner product on a vector space, but instead of vectors, we have functions. In this context, the inner product is often called the \textbf{integral inner product}. For two functions $f$ and $g$, the inner product is defined as:
    \[
    \inner{f}{g} = \int_0^L f(x) g(x) dx
    \]
    All of the same properties of an inner product hold for this integral inner product, such as linearity, symmetry, and positive-definiteness.
\end{definition}

\begin{definition}\label{def.orthogonality}
    Two functions $f$ and $g$ are said to be \textbf{orthogonal} with respect to the inner product if $\inner{f}{g} = 0$. A set of functions $\{f_n\}$ is said to be an \textbf{orthogonal set} if every pair of distinct functions in the set is orthogonal. If, in addition, each function in the set has norm 1 (i.e., $\inner{f_n}{f_n} = 1$), then the set is called an \textbf{orthonormal set}.
\end{definition}

\begin{lemma}\label{lemma.ortho_implies_lin_indep}
    An orthogonal set of functions is linearly independent.
\end{lemma}

With this dump of definitions and the lemma, we can now say that if we can show that the set of functions $\{\sin\left(\frac{n \pi x}{L}\right) \st n \in \NN \}$ is orthogonal with respect to the integral inner product, then we can conclude that they are linearly independent therefore they span some subspace of the function space. 

\subsubsection{Orthogonality Proof}\label{subsubsec.ortho_proof}
Lets check orthogonality of the sine functions by choosing two distinct functions from the set, say $\sin\left(\frac{m \pi x}{L}\right)$ and $\sin\left(\frac{n \pi x}{L}\right)$.
\[
\inner{\sin \tup{\frac{m \pi x}{L}}}{\sin \tup{\frac{n \pi x}{L}}} = \int_0^L \sin \tup{\frac{m \pi x}{L}} \sin \tup{\frac{n \pi x}{L}} dx
\]
\recall the trigonometric identity:
\[
\sin A \sin B = \frac{1}{2} \tup{\cos(A-B) - \cos(A+B)}
\]
\[
\implies \int_0^L \sin \tup{\frac{m \pi x}{L}} \sin \tup{\frac{n \pi x}{L}} dx = \frac{1}{2} \int_0^L \tup{\cos\left(\frac{(m-n) \pi x}{L}\right) - \cos\left(\frac{(m+n) \pi x}{L}\right)} dx
\]
This leads us to two cases: $m =n$, and $m \neq n$.\\
\textit{Case 1: $m = n$:}
\[
\frac{1}{2} \int_0^L \tup{\cos(0) - \cos\left(\frac{2n \pi x}{L}\right)} dx = \frac{1}{2} \int_0^L \tup{1 - \cos\left(\frac{2n \pi x}{L}\right)} dx
\]
\[
= \frac{1}{2} \left[ x - \frac{L}{2n \pi} \sin\left(\frac{2n \pi x}{L}\right) \right]_0^L = \frac{L}{2}
\]
which is not zero, so the functions are not orthogonal when $m = n$.\\
\textit{Case 2: $m \neq n$:}
\[
\frac{1}{2} \int_0^L \tup{\cos\left(\frac{(m-n) \pi x}{L}\right) - \cos\left(\frac{(m+n) \pi x}{L}\right)} dx
\]
\[
= \frac{1}{2} \left[ \frac{L}{(m-n) \pi} \sin\left(\frac{(m-n) \pi x}{L}\right) - \frac{L}{(m+n) \pi} \sin\left(\frac{(m+n) \pi x}{L}\right) \right]_0^L 
\]
Since $(m+n)$ and $(m-n)$ are both integers, the sine terms will evaluate to zero at both limits of integration, so the entire expression evaluates to zero. Thus,
\[
\inner{\sin \tup{\frac{m \pi x}{L}}}{\sin \tup{\frac{n \pi x}{L}}} = 0 \quad \text{for } m \neq n
\]

\subsection{Finding the Coefficients}\label{subsec.finding_coeffs}
Now that we have shown that the sine functions are orthogonal, we can use this property to find the coefficients $b_n$ in the Fourier series expansion of $f(x)$. We want to find $b_n$ such that:
\[
f(x) = \sum_{n=1}^\infty b_n \sin\left(\frac{n \pi x}{L}\right)
\]
To find $b_n$, we can take the inner product of both sides of the equation with $\sin\left(\frac{m \pi x}{L}\right)$ for some integer $m$:
\[
\inner{f}{\sin \tup{\frac{m \pi x}{L}}} = \inner{\sum_{n=1}^\infty b_n \sin\left(\frac{n \pi x}{L}\right)}{\sin \tup{\frac{m \pi x}{L}}}
\]
Using linearity of the inner product, we can pull the sum and the coefficients out:
\[
\inner{f}{\sin \tup{\frac{m \pi x}{L}}} = \sum_{n=1}^\infty b_n \inner{\sin\left(\frac{n \pi x}{L}\right)}{\sin \tup{\frac{m \pi x}{L}}}
\]
From the orthogonality proof, we know that $\inner{\sin\left(\frac{n \pi x}{L}\right)}{\sin \tup{\frac{m \pi x}{L}}} = 0$ for $n \neq m$ and $\inner{\sin\left(\frac{n \pi x}{L}\right)}{\sin \tup{\frac{n \pi x}{L}}} = \frac{L}{2}$ for $n = m$. Thus, the only term in the sum that survives is the one where $n = m$, giving us:
\[
\inner{f}{\sin \tup{\frac{m \pi x}{L}}} = b_m \inner{\sin\left(\frac{m \pi x}{L}\right)}{\sin \tup{\frac{m \pi x}{L}}} = b_m \frac{L}{2}
\]
Solving for $b_m$, we get:
\[
b_m = \frac{2}{L} \inner{f}{\sin \tup{\frac{m \pi x}{L}}} = \frac{2}{L} \int_0^L f(x) \sin\left(\frac{m \pi x}{L}\right) dx
\]
Thus, we have found the formula for the coefficients $b_n$ in the Fourier series expansion of $f(x)$:
\[
b_n = \frac{2}{L} \int_0^L f(x) \sin\left(\frac{n \pi x}{L}\right) dx, \quad n \in \NN
\]